\documentclass[../../../include/open-logic-section]{subfiles}

\begin{document}

\olfileid[id]{sth}{choice}{intro}

\olsection{Pengantar}

Dalam \crefrange{sth:cardinals::chap}{sth:card-arithmetic::chap}, kita
mengembangkan suatu teori kardinal dengan memperlakukan kardinal sebagai
ordinal. Pendekatan itu bergantung pada Aksioma Pengurutan Baik. Ternyata
Pengurutan Baik ekuivalen dengan prinsip lain---Aksioma Pilihan---dan telah
terjadi pembahasan filosofis yang serius mengenai dapat diterima atau
tidaknya prinsip tersebut. Pertanyaan kita dalam bab ini adalah: Bagaimana
Aksioma itu digunakan, dan dapatkah ia dibenarkan?

\end{document}
